\documentclass[11pt]{article}

\usepackage[a4paper,margin=1in]{geometry}
\usepackage{amsmath, amssymb}
\usepackage{enumitem}
\usepackage{hyperref}
\usepackage{setspace}

\setstretch{1.1}

\title{\textbf{CSE400 -- Fundamentals of Probability in Computing} \\[0.5em]
\large Lecture 10: Randomized Min-Cut Algorithm}
\author{}
\date{}

\begin{document}
\maketitle

\vspace{-1em}

\noindent
\textbf{Source:} Lecture 10 slides

\bigskip

\noindent
This document is a \textbf{strict, exam-ready lecture scribe for Lecture 10}, written \textbf{only from the provided lecture slides}.  
All definitions, notation, statements, algorithms, and explanations follow the \textbf{exact lecture flow and reasoning}, and it is intended to serve as a \textbf{sole reference for a closed-notes exam}.

%-------------------------------------------------

\section{Min-Cut Problem}

\subsection{Why Use Min-Cut?}

The min-cut problem is studied because it is useful for solving problems related to:

\begin{itemize}
    \item \textbf{Network connectivity}
    \item \textbf{Network reliability}
    \item \textbf{Optimization}
\end{itemize}

The lecture motivates min-cut through the following applications:

\begin{enumerate}
    \item \textbf{Network Design} \\
    Min-cut helps improve communication efficiency and optimize network flow.  
    By finding the \textbf{minimum capacity cut}, bottlenecks in the network can be identified.

    \item \textbf{Communication Networks} \\
    Min-cut is useful for understanding the \textbf{vulnerability of networks to failures}.  
    It helps in building \textbf{robust and fault-tolerant communication networks}.

    \item \textbf{VLSI Design} \\
    In \textbf{Very Large Scale Integration (VLSI)} design, min-cut algorithms help in \textbf{partitioning circuits} into smaller components, reducing \textbf{interconnectivity complexity}.
\end{enumerate}

%-------------------------------------------------

\section{What Is Min-Cut?}

\subsection{Cut-Set}

\textbf{Definition (Cut-Set):} \\
A \textbf{cut-set} in a graph is a set of edges whose removal breaks the graph into \textbf{two or more connected components}.

\medskip
\textbf{Reasoning:} \\
Edges in a cut-set are essential for connectivity. Removing them disconnects the graph, allowing us to study the weakest points of the network.

\subsection{Minimum Cut}

Given a graph
\[
G = (V, E)
\]
with $n$ vertices, the \textbf{minimum cut (min-cut) problem} is defined as:

\begin{quote}
To find a \textbf{minimum cardinality cut-set} in $G$.
\end{quote}

\noindent
\textbf{Explanation:} \\
Among all cut-sets, the one with the \textbf{smallest number of edges} is preferred, as it represents the smallest disconnection of the graph.

\subsection{Sensitivity of Min-Cut Algorithms}

Min-cut algorithms such as \textbf{Karger’s algorithm} are \textbf{randomized} and can be \textbf{sensitive to the initial choice of edges}.

\begin{itemize}
    \item Contracting \textbf{critical edges early} may destroy the true minimum cut.
    \item This explains why one execution may succeed while another may fail.
\end{itemize}

%-------------------------------------------------

\section{Edge Contraction}

\subsection{Main Operation}

The \textbf{main operation} in the min-cut algorithm is \textbf{edge contraction}.

\medskip
\textbf{Definition (Edge Contraction):} \\
Edge contraction removes an edge while \textbf{merging its two endpoint vertices} into one.

\subsection{Contracting an Edge}

When contracting an edge $(u, v)$:

\begin{enumerate}
    \item Vertices $u$ and $v$ are \textbf{merged into one vertex}.
    \item All edges between $u$ and $v$ are \textbf{eliminated}.
    \item All other edges are \textbf{retained}.
    \item The resulting graph:
    \begin{itemize}
        \item May contain \textbf{parallel edges}
        \item Contains \textbf{no self-loops}
    \end{itemize}
\end{enumerate}

\noindent
\textbf{Reasoning:} \\
Edge contraction reduces graph size while preserving information about possible cuts.

%-------------------------------------------------

\section{Successful and Unsuccessful Min-Cut Runs}

\subsection{Successful Min-Cut Run}

A \textbf{successful min-cut run} is one in which:

\begin{quote}
The algorithm correctly finds the \textbf{minimum cut} of the graph.
\end{quote}

\noindent
This occurs when edges belonging to the minimum cut are not contracted prematurely.

\subsection{Unsuccessful Min-Cut Run}

An \textbf{unsuccessful min-cut run} occurs when:

\begin{quote}
An edge in the true minimum cut is contracted early, leading to an incorrect result.
\end{quote}

\noindent
The structure of the minimum cut is destroyed due to early contraction.

%-------------------------------------------------

\section{Max-Flow Min-Cut Theorem}

\subsection{Theorem Statement}

\textbf{Max-Flow Min-Cut Theorem:}

\begin{quote}
In a flow network, the maximum flow from the source to the sink equals the total weight of the edges in a minimum cut.
\end{quote}

\subsection{Definitions}

\begin{itemize}
    \item \textbf{Capacity of a cut}: Sum of capacities of edges directed from set $X$ to set $Y$
    \item \textbf{Minimum cut}: Cut with the smallest capacity
    \item \textbf{Minimum cut capacity}: Capacity value of the minimum cut
    \item \textbf{Maximum flow}: Largest possible flow from source $S$ to sink $T$
\end{itemize}

\noindent
\textbf{Reasoning:} \\
The maximum achievable flow is exactly bounded by the smallest cut separating $S$ and $T$.

%-------------------------------------------------

\section{Deterministic Min-Cut Algorithm}

\subsection{Stoer--Wagner Min-Cut Algorithm}

Let $s$ and $t$ be vertices of a graph $G$.  
Let $G/\{s,t\}$ be the graph formed by merging $s$ and $t$.

\medskip
A minimum cut of $G$ is the smaller of:
\begin{itemize}
    \item A minimum $s$--$t$ cut of $G$
    \item A minimum cut of $G/\{s,t\}$
\end{itemize}

\subsection{Reasoning}

\begin{itemize}
    \item If a minimum cut separates $s$ and $t$, then the $s$--$t$ cut is optimal.
    \item Otherwise, merging $s$ and $t$ preserves the minimum cut.
\end{itemize}

\subsection*{Algorithm 1: \texttt{MinimumCutPhase$(G,a)$}}

\begin{enumerate}
    \item Initialize $A \leftarrow \{a\}$
    \item While $A \neq V$:
    \begin{itemize}
        \item Add the most tightly connected vertex to $A$
    \end{itemize}
    \item Return the cut weight
\end{enumerate}

\subsection*{Algorithm 2: \texttt{MinimumCut$(G)$}}

\begin{enumerate}
    \item While $|V| \ge 1$:
    \begin{itemize}
        \item Choose any vertex $a \in V$
        \item Run \texttt{MinimumCutPhase$(G,a)$}
        \item Store the cut if it is lighter
        \item Merge the last two added vertices
    \end{itemize}
    \item Return the minimum cut
\end{enumerate}

\noindent
\textbf{Property:} This algorithm always produces the \textbf{exact minimum cut}.

%-------------------------------------------------

\section{Randomized Min-Cut Algorithm}

\subsection{Why Randomized Algorithms?}

Randomized algorithms offer:
\begin{itemize}
    \item Probabilistic guarantees
    \item Improved efficiency
    \item Parallelizability
    \item Approximation guarantees
    \item Avoidance of worst-case scenarios
\end{itemize}

\subsection{Karger’s Randomized Algorithm}

\textbf{Idea:} Repeatedly select a random edge and contract it until only two vertices remain.

\begin{itemize}
    \item Remaining edges form a cut
    \item Random choices may or may not preserve the minimum cut
\end{itemize}

\subsection*{Algorithm 3: \texttt{RECURSIVE-RANDOMIZED-MIN-CUT$(G,\alpha)$}}

\textbf{Input:} Undirected multigraph $G$ with $n$ vertices, constant $\alpha > 0$ \\
\textbf{Output:} A cut $C$

\begin{enumerate}
    \item If $n \le 3$, find min-cut by exhaustive search
    \item Else:
    \begin{itemize}
        \item Repeat $\alpha$ times:
        \item Apply $n - \dfrac{n}{\sqrt{\alpha}}$ random contractions
        \item Recursively compute cut $C'$
        \item Keep the smallest cut
    \end{itemize}
    \item Return $C$
\end{enumerate}

%-------------------------------------------------

\section{Comparison: Deterministic vs Randomized Min-Cut}

\subsection*{Deterministic Min-Cut}

\begin{itemize}
    \item Always exact
    \item Higher time complexity
    \item Complexity:
    \[
    O(VE + V^2 \log V)
    \]
\end{itemize}

\subsection*{Randomized Min-Cut}

\begin{itemize}
    \item High-probability correctness
    \item Complexity:
    \[
    O(V^2)
    \]
\end{itemize}

%-------------------------------------------------

\section{Probability of Success}

The randomized min-cut algorithm outputs a minimum cut with probability at least:
\[
\frac{2}{n(n-1)}
\]

\noindent
Repeating the algorithm increases the probability of success.

%-------------------------------------------------

\section{Python Simulation}

Students are instructed to:
\begin{itemize}
    \item Open the \textbf{Campuswire post for Lecture 10}
    \item Download the provided \texttt{.ipynb} file
    \item Simulate the randomized min-cut algorithm
\end{itemize}

\bigskip
\noindent
\textbf{End of Lecture 10}

\end{document}
